% Autor: Elias Welt
\section{Einführung und Zielsetzung}
Das Frontend der Anwendung "`EduShelf"' repräsentiert die Schnittstelle zwischen dem Benutzer und den Backend-Diensten. Ziel der Implementierung war die Schaffung einer modernen, performanten und intuitiven Benutzeroberfläche, die das Lernen und Verwalten von Dokumenten so effizient wie möglich gestaltet.

Im Gegensatz zu klassischen Webanwendungen, die auf dem "`Multi-Page Application"' (MPA) Prinzip basieren und bei jeder Interaktion eine neue HTML-Seite vom Server anfordern, wurde EduShelf als "`Single Page Application"' (SPA) konzipiert. In einer SPA wird lediglich beim ersten Aufruf eine einzige HTML-Datei (\texttt{index.html}) geladen. Alle weiteren Interaktionen finden clientseitig statt, indem JavaScript (in diesem Fall React) den DOM (Document Object Model) dynamisch manipuliert.

Die Vorteile dieses Architekturansatzes sind vielfältig:
\begin{itemize}
    \item \textbf{Performance}: Da nach dem initialen Laden nur noch Rohdaten (im JSON-Format) und keine vollständigen HTML-Seiten mehr übertragen werden müssen, reduziert sich das Datenvolumen signifikant.
    \item \textbf{Benutzererfahrung (UX)}: Übergänge zwischen verschiedenen Ansichten erfolgen nahezu latenzfrei, ohne das typische "Flackern" eines Seiten-Reloads. Dies vermittelt das Gefühl einer nativen Desktop-Anwendung.
    \item \textbf{Entkopplung}: Durch die strikte Trennung von Frontend und Backend über eine REST-API können beide Systeme unabhängig voneinander entwickelt, getestet und skaliert werden.
\end{itemize}
